\section{Background}
\label{sec:background}

\subsection{Short-Burst (G.6)}
\label{subsec:background-sb}

Short-Burst~\citep{cannon2022} resolves the tension between mixing (long chains)
and optimisation (short chains) by running $n_{\mathrm{bursts}}$ short chains of
length $\ell_{\mathrm{burst}}$, keeping only the endpoints, and selecting the
best by edge-cut.
Each burst restarts from the previous best endpoint, producing a progressive
search that concentrates probability mass in the burst neighbourhood of each
successive plan.

In the G.6 implementation, the inner chain is \recom{}~\citep{deford2021}:
merge two adjacent districts, sample a random spanning tree, cut at the
balance point.
\recom{} accepts at $60$--$80\%$ on typical census-tract graphs.

\paragraph{Seeding.}
For burst $i$ with base seed $b$:
\[
  s^{(i)} =
  \mathrm{SHA256}\!\left(\texttt{"SB\_CHAIN\_"} \,\|\, i^{\mathrm{le64}}
    \,\|\, \texttt{"\_"} \,\|\, b^{\mathrm{le64}}\right)[0{:}8]\;\text{as } u64.
\]
Full reproducibility follows: any burst trajectory is recoverable from $b$
alone.

\subsection{Forest ReCom (G.9)}
\label{subsec:background-frecom}

\frecom{}~\citep{mccartan2020,g9paper} uses the same proposal as \recom{}
(merge two adjacent districts, sample a random spanning tree, cut at the
balance point) but adds a Metropolis--Hastings correction:
\[
  \alpha_{\mathrm{FR}}(\pi \to \pi') =
  \min\!\left(1,\; \frac{Q(\pi' \to \pi)}{Q(\pi \to \pi')}\right).
\]
Computing $Q(\pi' \to \pi)$ requires a \emph{reverse RNG} stream in addition
to the forward RNG used to sample the proposal.
The correction makes the stationary distribution approximately uniform over
valid plans; acceptance rate is $\approx 40$--$50\%$.

\subsection{Merge-Split (G.10)}
\label{subsec:background-ms}

\ms{}~\citep{janson2022,g10paper} applies the same merge-and-resample proposal
but computes the MH ratio via a single spanning-tree-count ratio rather than
the two-tree reverse evaluation used by \frecom{}.
This makes \ms{} cheaper per step while still providing a distributional
correction; acceptance rate is $\approx 55$--$65\%$.

\subsection{The Chain-Type Composition}
\label{subsec:background-composition}

All three chains---\recom{}, \frecom{}, and \ms{}---operate on the same
state space (valid $k$-district plans of the census-tract adjacency graph)
and propose moves via the same merge-and-resample operation.
They differ only in their acceptance criterion and the RNG resources required
per step.

This structural similarity makes the chain type a clean parameter for the
\sburst{} wrapper: swapping the inner chain from \recom{} to \frecom{} or
\ms{} changes the per-step cost and acceptance rate, but leaves the burst
structure, the endpoint-selection logic, and the seed architecture intact.
The G.12 implementations exploit this directly: \sbf{} and \sbms{} reuse
the \sburst{} burst loop, substituting the chain type.
